\section{Sprint 3 -- Hoàn thiện sản phẩm và bản sửa lỗi}\label{ch:sprint3}

Sau khi bàn giao \code{v1.0.0}, nhóm tự kiểm thử trên production và thấy một số điểm chưa đạt yêu cầu ``giống trang bán hàng thật'' của giảng viên:
\begin{itemize}
\item Giao diện chỉ cho chọn một topping, trong khi API cho phép nhiều.
\item Khách không biết trước còn đặt được bao nhiêu ly với lựa chọn của mình.
\item Nhân viên không thấy công thức từng món và số ly còn làm được.
\item US-20 chưa có test E2E trên trình duyệt.
\end{itemize}

Sprint 3 xử lý các điểm này mà không thay đổi tạo phẩm đã nộp của Task 1--5. Phần FR-03 được làm rõ trong SRS, kèm mục ``Lịch sử thay đổi''.

\textbf{Sprint Goal:} sản phẩm chỉn chu và đúng thực tế: nhiều topping, số ly theo công thức cho khách và nhân viên, test E2E trên trình duyệt.

{\footnotesize
\begin{longtable}{@{}L{1.1cm}L{4.3cm}L{1.7cm}L{1.6cm}L{0.9cm}L{1.7cm}L{2.6cm}@{}}
\caption{Product Backlog Sprint 3}\label{tab:backlogs3}\\
\toprule
\textbf{Mã} & \textbf{User story} & \textbf{Yêu cầu} & \textbf{MoSCoW} & \textbf{Point} & \textbf{Người làm} & \textbf{Trạng thái} \\ \midrule \endfirsthead
\toprule \textbf{Mã} & \textbf{User story} & \textbf{Yêu cầu} & \textbf{MoSCoW} & \textbf{Point} & \textbf{Người làm} & \textbf{Trạng thái} \\ \midrule \endhead
\bottomrule \endfoot
US-23 & Là khách hàng, tôi muốn chọn nhiều topping cho một ly và biết trước còn đặt được bao nhiêu ly & FR-03, FR-04, FR-16 & Should & 5 & A & Done (PR \#53) \\
US-24 & Là nhân viên, tôi muốn xem công thức từng món theo size và số ly còn làm được để chủ động nhập kho & FR-16, FR-17 & Should & 5 & A & Done (PR \#53) \\
US-25 & Là khách hàng, tôi muốn giao diện rõ ràng trên điện thoại: trang chủ, ảnh món, theo dõi từng đơn, cài như ứng dụng & NFR-05, NFR-06, FR-10, FR-11 & Should & 3 & A & Done (PR \#53) \\
\end{longtable}
}

\subsection{Nhiều topping và số ly làm được (US-23)}
Trang chi tiết món cho chọn tối đa 5 topping, giá cộng đủ từng topping.

Số ly tối đa của mỗi dòng giỏ hàng được API mới \code{POST /\allowbreak{}products/\allowbreak{}availability} tính theo:
\begin{itemize}
\item size, mức đường, mức đá, topping;
\item \emph{các dòng khác trong giỏ}: hai dòng dùng chung nguyên liệu thì tranh nhau phần còn lại.
\end{itemize}

Kết quả giới hạn 50 ly mỗi dòng và không trả số tồn kho thật. Khi còn dưới 50 ly, giao diện hiện ``Còn thêm được tối đa N ly''. Khi vượt, giao diện không cho thêm hoặc thanh toán.

API này dùng chung hàm tính với kiểm tra lúc \code{POST /orders}. Test INT-R06 khẳng định hai bên cho cùng kết quả: availability báo không đủ thì đặt đơn cũng bị từ chối, báo đủ thì đặt được.

Các tình huống số lượng được lấy từ Bài tập lớn 1 làm tiêu chí và viết thành integration test (INT-R01--R05, bảng \ref{tab:intr}):
\begin{itemize}
\item size L không đủ dù size S còn đủ;
\item topping hết nguyên liệu: ly có topping báo hết, ly không topping vẫn đặt được;
\item hai dòng tranh cùng nguyên liệu;
\item giới hạn 50 ly mỗi dòng.
\end{itemize}

\haianh{t5-32-chi-tiet-mon-nhieu-topping.jpg}{t5-33-gio-hang-23-ly-3-topping.jpg}{Chi tiết món: size, nhiều topping, mức đường, đá, số lượng}{Giỏ hàng: 23 ly Cà phê sữa đá size L với 3 topping, 62.000 đồng/ly}{Chọn nhiều topping và giỏ hàng}

\subsection{Công thức và số ly cho nhân viên (US-24)}
API \code{GET /\allowbreak{}ingredients/\allowbreak{}recipes} (chỉ STAFF) trả cho từng món:
\begin{itemize}
\item lượng nguyên liệu cho 1 ly theo size S/M/L;
\item số ly S/M/L làm được với tồn kho hiện tại;
\item nguyên liệu giới hạn.
\end{itemize}

Trang \code{/\allowbreak{}staff/\allowbreak{}recipes} hiển thị bảng công thức và cho nhập kho ngay tại chỗ. Nhân viên được xem số thật, nên số ly không bị giới hạn ở 50. Trang kho \code{/\allowbreak{}staff/\allowbreak{}ingredients} có thêm cột ``Dùng cho'': mỗi nguyên liệu dùng cho món nào, còn bao nhiêu ly.

\textbf{Kiểm thử trên production} (hình \ref{fig:recipebefore}): nhập thêm 200 g đá viên. Tồn kho từ 13.960 g lên 14.160 g, số ly Cà phê sữa đá cập nhật ngay từ 116/93/77 lên 118/94/78.

Tính lại theo công thức (đá viên là nguyên liệu giới hạn, 120 g cho 1 ly size S):
\begin{itemize}
\item Size S: $14.160 / 120 = 118$.
\item Size M (150 g/ly): $\lfloor 14.160 / 150 \rfloor = 94$.
\item Size L (180 g/ly): $\lfloor 14.160 / 180 \rfloor = 78$.
\end{itemize}

\haianh{t5-59-cong-thuc-so-ly-truoc-nhap.jpg}{t5-60-cong-thuc-so-ly-sau-nhap-200g-da.jpg}{Trước khi nhập: Cà phê sữa đá 116/93/77 ly, giới hạn bởi đá viên}{Sau khi nhập 200 g đá viên: 118/94/78 ly}{Trang công thức và số ly trước và sau khi nhập kho\label{fig:recipebefore}\label{fig:recipeafter}}

\subsection{Giao diện, theo dõi đơn và PWA (US-25)}
\begin{itemize}
\item Trang chủ mới.
\item Hình minh họa đồ uống bằng SVG tự vẽ, thay cho ảnh món bị lỗi.
\item Trang theo dõi đơn \code{/orders/[id]}: thanh tiến trình 5 bước (Đã đặt, Đã thanh toán, Đang pha chế, Sẵn sàng nhận, Hoàn tất), tự cập nhật.
\item Manifest PWA.
\item Liên kết ``Quầy nhân viên'' cho tài khoản STAFF.
\item Không cuộn ngang ở màn hình 375 px.
\end{itemize}

\anh{t5-56-khach-theo-doi-don-2-hoan-tat.jpg}{0.9\linewidth}{Trang theo dõi đơn \#2 của khách: tiến trình ở bước Hoàn tất, thanh toán VNPay sandbox thành công}

\subsection{Test E2E Playwright trong CI (US-20)}
Thư mục \code{tests/e2e} có 7 kịch bản (8 test) chạy trên trình duyệt thật. Cả API, web và PostgreSQL đều thật, database đã migrate và seed (mục \ref{sec:e2e}).

CI có thêm job thứ hai ``E2E (Playwright)'', thực hiện lần lượt:
\begin{enumerate}
\item Khởi động PostgreSQL riêng.
\item Migrate và seed.
\item Build.
\item Chạy API và web.
\item Chạy Playwright trên Chromium.
\end{enumerate}

Khi lỗi, báo cáo và log được lưu thành artifact. Job chạy trên \code{ubuntu-24.04} vì Playwright 1.56 chưa hỗ trợ Ubuntu 26.04 (runner \code{ubuntu-latest} hiện tại). Đây là một trong hai lỗi cấu hình phát hiện ngay trên CI của PR \#53. Lỗi còn lại là gọi sai tên script seed (\code{seed} thay vì \code{db:seed}).

\haianh{t5-27-ci-pr53-2-job-e2e-playwright.jpg}{t5-28-cd-v1.1.0-success.jpg}{CI của PR \#53: hai job đều thành công}{CD tag \code{v1.1.0}: CI 2 job $\rightarrow$ GHCR $\rightarrow$ Render và smoke test}{CI hai job và CD \code{v1.1.0}}

\subsection{Phát hành \code{v1.1.0}}
Tag \code{v1.1.0} (commit \code{2eb672f}) chạy CD gồm cả job E2E, và xanh. Production trả \code{"version":"v1.\allowbreak{}1.\allowbreak{}0"}. Bản này không có migration mới và không đổi biến môi trường.


\subsection{Kiểm thử thủ công trên production và bản sửa lỗi \code{v1.1.1}}
Ngày 17/11/2026, nhóm kiểm thử thủ công toàn bộ luồng trên production \code{v1.1.0} và chụp 41 ảnh bằng chứng (Phụ lục \ref{app:bangchung}). Lần kiểm thử này phát hiện 3 lỗi:

\begin{longtable}{@{}L{3.6cm}L{4.4cm}L{5.8cm}@{}}
\caption{Lỗi phát hiện khi kiểm thử production và cách sửa ở \code{v1.1.1}}\label{tab:loi111}\\
\toprule
\textbf{Lỗi} & \textbf{Nguyên nhân} & \textbf{Cách sửa và kiểm chứng} \\ \midrule \endfirsthead
\toprule \textbf{Lỗi} & \textbf{Nguyên nhân} & \textbf{Cách sửa và kiểm chứng} \\ \midrule \endhead
\bottomrule \endfoot
Nhân viên bấm ``Đăng xuất'' thì ở lại form đăng nhập nhân viên, không về trang web & Nút chỉ xóa phiên và đổi trạng thái của trang \code{/staff} & Nút là liên kết về \code{/}, xóa token và thông tin người dùng trước khi chuyển trang. Test \code{staff/\allowbreak{}staff.\allowbreak{}test}; ảnh 66, 67 \\
Chữ ``100\%'' ở nút mức đường, mức đá bị cắt trên màn hình rộng & Lớp \code{px-4} mặc định đè lên \code{px-0} & Hàm tạo lớp nút nhận khoảng đệm riêng; nút không xuống dòng; hai nhóm nằm cạnh nhau từ màn hình \code{xl}. Kiểm tra ở 1920 px và 1024 px \\
Chrome báo không tải được biểu tượng ứng dụng trong manifest & Manifest chỉ có biểu tượng SVG; image web không chép thư mục \code{public} & Thêm \code{icon-192.png}, \code{icon-512.png}; Dockerfile chép \code{public}. Test \code{page.test}; ảnh 63 (trước), 65 (sau) \\
\end{longtable}

\textbf{Vì sao E2E không bắt được lỗi biểu tượng.} E2E-T06b chỉ kiểm tra biểu tượng đầu tiên trong manifest trả về được qua HTTP. Test chạy bằng \code{npm run start}, không chạy trên image Docker. Nó không kiểm tra Chrome có dùng được biểu tượng để cài ứng dụng hay không.

Hai lỗi giao diện còn lại nằm ngoài các kịch bản E2E hiện có. Nhóm ghi các điểm này vào Retrospective Sprint 3.

PR \#55 sửa cả ba lỗi. Tag \code{v1.1.1} (commit \code{6229753}) chạy CD xanh trong 10 phút 56 giây. Release notes: \code{docs/\allowbreak{}release/\allowbreak{}release-notes-v1.\allowbreak{}1.\allowbreak{}1.\allowbreak{}md}.

\anh{t5-64-cd-v1.1.1-success.jpg}{0.9\linewidth}{CD tag \code{v1.1.1}: CI 2 job, build và push GHCR 7 phút 8 giây, deploy Render và smoke test 1 phút 22 giây}

\subsection{Sprint Review và Retrospective Sprint 3}
\textbf{Kết quả:} đạt Sprint Goal, phát hành \code{v1.1.0} và \code{v1.1.1} cùng ngày 17/11/2026.

\textbf{Điều làm tốt:}
\begin{itemize}
\item Lấy tình huống số lượng của Bài tập lớn 1 làm tiêu chí, viết test ở cả mức unit và integration trước khi làm giao diện.
\item API số ly và kiểm tra lúc đặt đơn dùng chung một hàm tính.
\item Job E2E bắt lỗi cấu hình ngay trong PR.
\end{itemize}

\textbf{Điều chưa tốt:}
\begin{itemize}
\item Toàn bộ Sprint 3 do một người hiện thực; các thành viên khác chưa tham gia review mã.
\item Hai lần CI đỏ vì cấu hình.
\item Test số lượng của Bài tập lớn 1 lẽ ra nên chuyển sang từ Sprint 2.
\item Ba lỗi giao diện của \code{v1.1.0} chỉ lộ ra khi thử tay trên production.
\end{itemize}

\textbf{Hành động:}
\begin{itemize}
\item Mỗi PR có ít nhất một thành viên khác review và approve.
\item Cố định phiên bản hệ điều hành của runner cho job phụ thuộc trình duyệt.
\item Chạy thử job CI mới trên nhánh nháp trước khi gắn vào PR chính.
\item Thêm E2E cho đăng xuất nhân viên và cho manifest có biểu tượng PNG 192/512.
\item Xem giao diện ở 375 px và 1920 px trước khi phát hành.
\end{itemize}
